% ----------------------------------------------------------------------
%  Итоги главы 3: сводка, ответы, задачи, литература
% ----------------------------------------------------------------------
\clearpage
\section*{Итоги главы}
\addcontentsline{toc}{section}{Итоги главы}
\markright{Итоги главы 3}

Глава прошла путь от электронного строения атома до свойств кристалла.
Почти все свойства твёрдых тел определяются двумя «газами»~—
\emph{электронным} (электроны проводимости, подчиняющиеся статистике
Ферми) и \emph{фононным} (кванты колебаний решётки)~— и их рассеянием на
дефектах. Таблица~\ref{tab:ch3-summary} сводит вместе, что за какое
свойство отвечает.

\medskip\noindent\begin{minipage}{\textwidth}
\centering
\captionof{table}{Свойства твёрдых тел: носители, главные законы, чувствительность
к структуре}
\label{tab:ch3-summary}
\footnotesize
\begin{tabularx}{\textwidth}{@{}>{\raggedright\arraybackslash\hspace{0pt}}p{24mm}>{\raggedright\arraybackslash\hspace{0pt}}p{30mm}>{\raggedright\arraybackslash\hspace{0pt}}X>{\raggedright\arraybackslash\hspace{0pt}}p{24mm}@{}}
\toprule
Свойство & Кто отвечает & Главные закономерности & Структурная чувствительность \\
\midrule
Энергия связи, структура & валентные электроны & тип связи (ионная, ковалентная, металлическая, ван-дер-ваальсова); правило $8-N$; отношение радиусов & --- \\
Фазы в сплавах & электронная концентрация, размеры атомов & правила Юм-Розери; электронные соединения 3/2, 21/13, 7/4; фазы Лавеса, внедрения & --- \\
Диффузия & вакансии, межузельные атомы & законы Фика; $D=D_0e^{-Q/kT}$; эффект Киркендалла; диффузия по границам & сильная \\
Электропроводность & электроны у $\EF$ & $\sigma=ne^2l/mv_\text{F}$; $\rho\propto T$ при $T>\theta_D$; правила Маттиссена, Нордгейма, Курнакова & сильная (остаточное $\rho$) \\
Магнетизм & спины и орбиты электронов & Кюри, Кюри--Вейсс, Паули; обменный интеграл (Бете--Слэтер); домены & $I_s$~— нет; $H_\text{к}$, $\mu$~— сильная \\
Теплоёмкость & фононы; электроны & $3R$; $T^3$ (Дебай); $\gamma T$; $\lambda$-аномалии при переходах 2-го рода & нет \\
Теплопроводность & фононы; электроны & $\lambda=\frac13cvl$; $\lambda/\sigma=LT$ & сильная при низких $T$ \\
Тепловое расширение & ангармонизм колебаний & $\alpha\propto C_V$ (Грюнайзен); инвар & слабая \\
Упругость & силы связи & $E(T_\text{пл})$; $E(Z)$ периодична; $\Delta E$-эффект & слабая \\
Внутреннее трение & дефекты, релаксация & пик при $\omega\tau=1$; Снук, Зинер & очень сильная \\
\bottomrule
\end{tabularx}
\end{minipage}

\begin{itogi}[Главные формулы главы]
\begin{itemize}
\item Энергия Ферми: $\EF=\frac{\hbar^2}{2m}(3\pi^2n)^{2/3}$; плотность состояний
$g(E)\propto\sqrt E$; функция Блоха $\psi_{\vb k}=u_{\vb k}(\vb r)e^{i\vb k\vb r}$.
\item Ионный кристалл: $U=-\frac{\alpha e^2}{4\pi\varepsilon_0r}+\frac B{r^n}$,
$U_0=-\frac{\alpha e^2}{4\pi\varepsilon_0r_0}\bigl(1-\frac1n\bigr)$.
\item Диффузия: $J=-D\,\partial c/\partial x$, $\partial c/\partial t=D\,\partial^2c/\partial x^2$;
$\langle X^2\rangle=2Dt$; $D=f\gamma\lambda^2\nu$; Даркен: $\tilde D=c_BD_A+c_AD_B$.
\item Электропроводность: $\sigma=ne^2l_\text{F}/(mv_\text{F})$;
$\rho=\rho_\text{ост}+\rho_\text{фон}(T)$; Нордгейм $\Delta\rho\propto c(1-c)$; Холл $R=1/(ne)$.
\item Магнетизм: $\vb B=\mu_0(\vb H+\vb I)$; $\chi=C/(T-\theta)$; $M_\text{Б}=e\hbar/2m$.
\item Средняя энергия колебания $\hbar\omega/(e^{\hbar\omega/kT}-1)$;
$C_V=\frac{12\pi^4}{5}R(T/\theta_D)^3$; $C_\text{эл}=\frac{\pi^2}2R\,T/T_\text{F}$.
\item Теплопроводность: $\lambda=\frac13c\upsilon l$; $L=\frac{\pi^2}3(k/e)^2=2{,}45\cdot10^{-8}$~Вт$\cdot$Ом/К$^2$.
\item Расширение: $\alpha=\gamma_\text{Г}KC_V/3V$.
\item Упругость: $G=E/2(1+\nu)$, $D=E/3(1-2\nu)$; $E=4\rho\nu_\text{рез}^2L^2$.
\item Внутреннее трение: $Q^{-1}=\frac1{2\pi}\frac{\Delta W}W=\Delta\frac{\omega\tau}{1+\omega^2\tau^2}$;
$\ln\frac{\omega_2}{\omega_1}=\frac HR\bigl(\frac1{T_1}-\frac1{T_2}\bigr)$.
\end{itemize}
\end{itogi}

%=======================================================================
\section*{Ответы на вопросы для самопроверки}
\addcontentsline{toc}{section}{Ответы на вопросы для самопроверки}
\markright{Ответы на вопросы главы 3}

{\small
\paragraph{3.1.2 (строение атома)}
\begin{enumerate}[label=\arabic*.]
\item $\lambda=h/\sqrt{2mE}=6{,}63\cdot10^{-34}/\sqrt{2\cdot1{,}675\cdot10^{-27}\cdot0{,}025\cdot1{,}6\cdot10^{-19}}\approx1{,}8\cdot10^{-10}$~м~—
порядка межатомного расстояния; поэтому тепловые нейтроны дифрагируют
на кристаллах.
\item Слой $M$ ($n=3$)~— $2n^2=18$ электронов; оболочка $4f$ ($l=3$)~—
$2(2l+1)=14$.
\item Второй электрон приходится вырывать из заполненной оболочки
$1s^2$, которая находится гораздо ближе к ядру и почти не экранирована:
75,6~эВ против 5,4~эВ.
\item $4s\to3d\to4p$: по правилу $n+l$ у $4s$ сумма 4, у $3d$ и $4p$~— 5
(раньше заполняется меньшее $n$). Физически $4s$-электрон сильнее
проникает к ядру и хуже экранирован.
\item Между ними заполняется $4f$-оболочка (лантаниды); $f$-электроны
плохо экранируют заряд ядра, и радиусы следующих элементов «сжимаются»
(лантаноидное сжатие); у золота добавляется релятивистское сжатие
$6s$-орбитали.
\end{enumerate}

\paragraph{3.1.5 (типы связи)}
\begin{enumerate}[label=\arabic*.]
\item Натрий: у него электроотрицательность меньше (2,8 против 8,3~эВ),
электрон слабее связан с атомом и переходит к хлору.
\item Для NaCl правило даёт 1,61~эВ вместо опытных 4,26~эВ: связь не
ковалентная, а ионная, и большая часть энергии связи~— электростатическое
притяжение ионов.
\item Четыре шара радиуса $R$ в вершинах тетраэдра с ребром $2R$; расстояние
от центра до вершины $2R\sqrt{3/8}=R\sqrt{3/2}$; радиус пустоты
$r=R(\sqrt{3/2}-1)=0{,}225R$.
\item $r_+/r_-=0{,}078/0{,}220=0{,}355$~— тетраэдрический интервал
0,225--0,414, т.\,е. ожидается координация 4. На деле LiI имеет структуру
NaCl (координация 6): правило отношения радиусов~— лишь грубое
геометрическое приближение, особенно для очень малого катиона.
\item У RbBr $r_+/r_-=0{,}76$ лежит у границы 0,732; давление стабилизирует
более плотную фазу, а структура CsCl (координация 8) плотнее, чем NaCl.
\item С ростом числа электронов растёт поляризуемость атомов, а с ней~—
силы Ван-дер-Ваальса.
\end{enumerate}

\paragraph{3.2.1 (свободные электроны)}
\begin{enumerate}[label=\arabic*.]
\item Интервал между уровнями $\sim\EF/N\sim10^{-22}$~эВ для образца в
1~см$^3$~— несравнимо меньше $kT$ и любой энергии, которую электрон может
получить.
\item $\EF\propto n^{2/3}$: в $8^{2/3}=4$ раза.
\item $n=8{,}4\cdot10^{28}$~м$^{-3}$; $\EF=\frac{\hbar^2}{2m}(3\pi^2n)^{2/3}\approx1{,}1\cdot10^{-18}$~Дж~$\approx7$~эВ.
\item По принципу Паули электроны занимают все уровни вплоть до $\EF$;
средняя энергия $\frac35\EF\approx3$--$4$~эВ даже при 0~К.
\item Доля $\sim kT/\EF=0{,}026/3{,}2\approx1\,\%$.
\end{enumerate}

\paragraph{3.2.2 (зонная теория)}
\begin{enumerate}[label=\arabic*.]
\item Обе~— бегущие волны $e^{i\vb k\vb r}$, но у функции Блоха амплитуда
$u_{\vb k}(\vb r)$ модулирована с периодом решётки.
\item При $k=\pi/a$ выполняется условие Брэгга; отражённая и падающая
волны образуют две стоячие: $\cos^2$ собирает электроны на ионах (ниже
энергия), $\sin^2$~— между ними (выше энергия); разница~— щель.
\item Грани первой зоны ОЦК~— плоскости $\{110\}$; расстояние
$G_{110}/2=\pi\sqrt2/a\approx1{,}03\cdot10^{10}$~м$^{-1}$.
\item Зоны перекрываются: электроны начинают заполнять вторую зону
раньше, чем заполнится первая; обе зоны частично заполнены.
\item Каждому состоянию $\vb k$ соответствует состояние $-\vb k$ с
противоположной скоростью; поле не может изменить заполнение
полностью занятой зоны, и суммарный ток равен нулю.
\end{enumerate}

\paragraph{3.3 (фазы в сплавах)}
\begin{enumerate}[label=\arabic*.]
\item В растворе вычитания часть узлов одной из подрешёток пуста.
Различают по числу атомов в ячейке $n=\rho V_\text{яч}/(m_uA_\text{ср})$:
$n<n_0$~— вычитания, $n>n_0$~— внедрения.
\item Огромная разница атомных радиусов (около 50\,\%, далеко за пределом
15\,\%).
\item $e/a=1+2x=1{,}4$, $x=0{,}2$~— 20~ат.\,\% Al (опыт~— около 19,6\,\%).
\item Интенсивность сверхструктурных линий $\propto(f_\text{Ni}-f_\text{Fe})^2$, а
атомные номера 28 и 26 почти равны.
\item $(9\cdot1+4\cdot3)/13=21/13$~— $\gamma$-фаза (тип $\gamma$-латуни).
\item Октаэдрическая пустота ГЦК крупнее ($0{,}414R$ против $0{,}225R$).
\end{enumerate}

\paragraph{3.4 (диффузия)}
\begin{enumerate}[label=\arabic*.]
\item Движущая сила диффузии~— градиент химического потенциала, а не
концентрации; $D=MkT(1+\d\ln\gamma/\d\ln c)$, и термодинамический множитель
зависит от состава.
\item Толщина слоя $\propto\sqrt{Dt}$: в 9 раз.
\item Углерод~— атом внедрения: ему не нужна вакансия, и энергия
активации меньше (около 0,8~эВ против ${\approx}2{,}5$~эВ у самодиффузии
железа).
\item $f\approx1-2/z=0{,}75$; точно~— 0,727.
\item При граничной диффузии $\ln C$ на большой глубине линейно зависит от
глубины $y$ (у объёмной~— от $y^2$), а наклон убывает как $t^{-1/4}$.
\item Цинк диффундирует быстрее меди: из латуни уходит больше атомов, чем
приходит, и её объём уменьшается~— метки смещаются в её сторону.
\item Внутри спинодали $\partial^2G/\partial c^2<0$: термодинамический
множитель отрицателен, $\tilde D<0$, и поток направлен против градиента
концентрации.
\end{enumerate}

\paragraph{3.5 (электрические свойства)}
\begin{enumerate}[label=\arabic*.]
\item У Друде длина пробега постоянна, а скорость электронов~—
тепловая, $\propto\sqrt T$, откуда $\rho\propto\sqrt T$. На деле скорость~—
фермиевская (от $T$ не зависит), а длина пробега $\propto1/n_\text{ф}\propto1/T$.
\item Изменить состояние под действием поля могут только электроны у
поверхности Ферми: ниже все состояния заняты.
\item У металла проводимость падает с ростом $T$ (растёт рассеяние на
фононах), у полупроводника растёт экспоненциально (растёт число носителей).
\item Остаточное сопротивление при 4,2~К определяется примесями и
дефектами: чем больше отношение, тем чище металл.
\item Проводимость бериллия определяется в основном дырками (электроны у
верха почти заполненной зоны), и постоянная Холла положительна.
\item При 20\,°C образуются зоны Гинье--Престона размером порядка длины
волны электрона~— сильные рассеиватели; при 225\,°C выделяются крупные
частицы $\theta'$, а матрица обедняется медью.
\item В сверхпроводниках II рода вихри под действием силы Лоренца
движутся, что вызывает потери; пиннинг на дефектах удерживает вихри и
повышает критический ток.
\end{enumerate}

\paragraph{3.6 (магнитные свойства)}
\begin{enumerate}[label=\arabic*.]
\item Ларморова прецессия есть в любом атоме, но эффект мал и маскируется
парамагнетизмом, если у атомов есть собственные моменты; в чистом виде
он виден у веществ с заполненными оболочками.
\item Перевернуть спин по полю могут только электроны в узкой полосе у
$\EF$; их число задаётся плотностью состояний $g(\EF)$ и от температуры
не зависит. Моменты же ионов соли ориентируются полем против теплового
беспорядка, отсюда $\chi\propto1/T$.
\item Из 10 электронов в $3d$ и $4s$ около 0,6 находятся в $4s$-зоне, и в
$3d$-зоне остаётся 0,6 дырки~— некомпенсированных спинов.
\item У марганца отношение межатомного расстояния к диаметру
$3d$-оболочки мало, обменный интеграл отрицателен (кривая Бете--Слэтера);
в сплавах Гейслера атомы Mn раздвинуты, и $A>0$.
\item Домены уменьшают магнитостатическую энергию. В малой частице
энергия доменной стенки больше выигрыша от разбиения.
\item Дислокации и внутренние напряжения тормозят смещение доменных
стенок.
\item Цинк добавляет по 2 электрона на атом; 0,6 дырки заполнятся при
$c=0{,}6/2=30$~ат.\,\%.
\end{enumerate}

\paragraph{3.7 (тепловые свойства)}
\begin{enumerate}[label=\arabic*.]
\item Смысл имеют лишь смещения атомов, а волна короче $2a$ задаёт те же
смещения, что и некоторая волна длиннее $2a$.
\item Акустические~— соседние атомы движутся в фазе, $\omega\to0$ при
$k\to0$; оптические~— в противофазе, частота конечна при $k=0$. У кремния
есть: в примитивной ячейке два атома, значит три оптические ветви.
\item У Эйнштейна все колебания имеют одну частоту и вымерзают
одновременно, экспоненциально. У Дебая есть сколь угодно низкие частоты
(длинные звуковые волны), их число $\propto T^3$.
\item $T=\sqrt{\gamma/A}$, $A=12\pi^4R/5\theta_D^3=0{,}048$~мДж/(моль$\cdot$К$^4$);
$T\approx3{,}8$~К.
\item $\gamma\propto g(\EF)$, а у переходных металлов $\EF$ лежит в узкой
$d$-зоне с высокой плотностью состояний.
\item Магнитное~— второго рода (конечный $\lambda$-пик), $\alpha\to\gamma$~—
первого (скрытая теплота, бесконечно узкий пик).
\item Растёт число фононов, но падает их длина пробега. При низких $T$
пробег ограничен размером образца ($\lambda\propto LT^3$), поэтому чем больше
кристалл, тем выше максимум и тем ниже его температура.
\item Константан~— концентрированный неупорядоченный раствор: электроны
сильно рассеиваются на атомах другого сорта, и $\lambda_\text{эл}$ падает до
уровня решёточной.
\item Гармоническая яма симметрична: среднее положение атома не зависит
от амплитуды колебаний.
\end{enumerate}

\paragraph{3.8 (упругие свойства)}
\begin{enumerate}[label=\arabic*.]
\item При $\nu>0{,}5$ растяжение уменьшало бы объём, а модуль
$D=E/3(1-2\nu)$ стал бы отрицательным~— тело было бы неустойчиво.
\item Из-за ангармонизма решётка расширяется, межатомные силы на больших
расстояниях слабее; $e\propto\beta$, $e\approx\alpha/0{,}04$.
\item $\nu=\sqrt{E/\rho}/2L=5170/0{,}4\approx13$~кГц.
\item Модуль определяется силами связи во всём объёме, а границы
занимают ничтожную его долю; внутреннее трение создают именно дефекты.
\item Напряжение смещает доменные стенки и поворачивает намагниченность,
добавляя магнитострикционную деформацию; в насыщении этих процессов нет.
\item При быстром растяжении образец охлаждается и потому сокращается,
частично компенсируя удлинение,~— он «жёстче».
\item При $\omega\tau\gg1$ релаксация не успевает пойти, при $\omega\tau\ll1$
успевает полностью; в обоих случаях деформация идёт в фазе с напряжением.
\item $1/T_2=1/T_1-R\ln100/H$; $T_2\approx361$~К.
\end{enumerate}
}

%=======================================================================
\section*{Задачи}
\addcontentsline{toc}{section}{Задачи}
\markright{Задачи к главе 3}

\begin{enumerate}[label={\sffamily\bfseries\color{fmHead}\arabic*.}]
\item Найдите энергию и скорость Ферми алюминия ($\rho=2{,}70$~г/см$^3$,
$A=27$, три свободных электрона на атом).\\
\emph{Ответ:} $n=1{,}81\cdot10^{29}$~м$^{-3}$, $\EF\approx11{,}7$~эВ,
$v_\text{F}=\hbar k_\text{F}/m\approx2{,}0\cdot10^6$~м/с.
\item Сравните электронную теплоёмкость серебра ($\EF=5{,}5$~эВ) при 300~К
по классической теории Друде и по теории Зоммерфельда.\\
\emph{Ответ:} $\frac32R=12{,}5$~Дж/(моль$\cdot$К) против
$\frac{\pi^2}{2}R\,T/T_\text{F}\approx0{,}19$~Дж/(моль$\cdot$К) ($T_\text{F}=64\,000$~К).
\item При каком составе латуни сфера Ферми касается границ первой зоны
Бриллюэна ГЦК ($e/a=1{,}36$) и ОЦК ($e/a=1{,}48$) решёток?\\
\emph{Ответ:} $e/a=1+x$: 36 и 48~ат.\,\% Zn~— близко к пределам
$\alpha$- и $\beta$-фаз.
\item Выразите показатель $n$ отталкивания Борна через сжимаемость $K$
кристалла типа NaCl и оцените его ($r_0=0{,}282$~нм, $K=4{,}2\cdot10^{-11}$~Па$^{-1}$,
$\alpha=1{,}748$).\\
\emph{Ответ:} $n=1+\dfrac{18\cdot4\pi\varepsilon_0r_0^4}{\alpha e^2K}\approx8$.
\item На какой глубине у профиля диффузии из тонкого слоя вторая
производная концентрации обращается в нуль? Сравните со
среднеквадратичным смещением.\\
\emph{Ответ:} $x=\sqrt{2Dt}$~— ровно среднеквадратичное смещение атома
вдоль оси.
\item Найдите термодинамический множитель для регулярного раствора с
энергией смешения $\Omega$ и условие границы спинодали.\\
\emph{Ответ:} $1-2\Omega x(1-x)/kT$; спинодаль: $kT=2\Omega x(1-x)$.
\item Сколько независимых коэффициентов Онзагера описывают
изотермическую диффузию в $n$-компонентной системе при $\sum J_i=0$?\\
\emph{Ответ:} независимых потоков $n-1$, коэффициентов $(n-1)^2$, из них с
учётом соотношений взаимности независимы $n(n-1)/2$.
\item Оцените теплопроводность алюминия при 293~К по его удельному
сопротивлению $2{,}65\cdot10^{-8}$~Ом$\cdot$м.\\
\emph{Ответ:} $\lambda=LT/\rho\approx270$~Вт/(м$\cdot$К) (опыт~— 237).
\item Оцените температуру Дебая алюминия по скоростям звука
$\upsilon_l=6420$~м/с, $\upsilon_t=3040$~м/с ($\rho=2700$~кг/м$^3$).\\
\emph{Ответ:} $\upsilon_m\approx3420$~м/с, $\theta_D\approx400$~К
(калориметрия~— 428~К).
\item Амплитуда крутильных колебаний уменьшилась вдвое за 50 периодов.
Найдите внутреннее трение.\\
\emph{Ответ:} $Q^{-1}=\ln2/(50\pi)\approx4{,}4\cdot10^{-3}$.
\item Найдите постоянную Холла натрия ($n=2{,}65\cdot10^{28}$~м$^{-3}$).\\
\emph{Ответ:} $R=-1/(ne)\approx-2{,}4\cdot10^{-10}$~м$^3$/Кл.
\end{enumerate}

\paragraph{Вопросы для повторения (по исходному учебнику)}
{\small
\begin{enumerate}[label=\arabic*.]
\item Какова симметрия электронного облака заполненных $s$-, $p$-, $d$-,
$f$-подоболочек?
\item Как выглядела бы таблица Менделеева, если бы уровни заполнялись по
схеме одноэлектронного атома?
\item Рассчитайте межъядерное расстояние и энергию связи молекулы H$_2$.
\item Получите выражение для энергии связи металла в рамках ионной
модели.
\item Найдите постоянную Маделунга для кристалла типа CsCl.
\item Укажите классический предел статистики Ферми--Дирака.
\item Найдите плоскости, ограничивающие первую зону Бриллюэна ГПУ-решётки.
Какой вид имеет эта зона?
\item Покажите, что в одномерной решётке разрывы энергии на границе зоны
эквивалентны условию брэгговского отражения.
\item Получите радиус сферы Ферми, если $N$ атомов с $Z$ валентными
электронами занимают куб с ребром $L$.
\item Укажите знак активационного объёма диффузии для кольцевого,
прямого межузельного и вакансионного механизмов.
\item Как квантовая теория классифицирует твёрдые тела по электрическим
свойствам? Чем полупроводник отличается от диэлектрика?
\item Какими механизмами обусловлена температурная зависимость
электропроводности проводников и полупроводников?
\item Как классифицируются сверхпроводники?
\item Какими магнитными свойствами и почему обладают Bi, CaCl$_2$, Zr, Hg,
Na, Al, жидкий N$_2$?
\item Какие магнитные характеристики структурно чувствительны? Когда
наблюдается прямоугольная петля гистерезиса?
\item Какими характеристиками должны обладать ферромагнетики для
сердечников трансформаторов, экранов, постоянных магнитов, носителей
информации?
\item Почему теплоёмкость тел со сложной решёткой хуже описывается теорией
Дебая?
\item Почему алмаз при комнатной температуре имеет теплоёмкость заметно
меньше $3R$ и очень высокую теплопроводность?
\item Какие вклады в коэффициент расширения существуют помимо решёточного
и когда они существенны?
\item Докажите, что с уменьшением чистоты металла максимум его
низкотемпературной теплопроводности смещается к более высоким
температурам и понижается.
\item Докажите, что максимум внутреннего трения наблюдается при
$\omega\tau=1$.
\item Как зависят от температуры модули упругости и коэффициент Пуассона
чистых металлов?
\item При каких частотах наблюдаются пики внутреннего трения для разных
механизмов?
\item Как методом внутреннего трения определить энергию активации
диффузии примесей внедрения?
\end{enumerate}
}

%=======================================================================
\section*{Литература}
\addcontentsline{toc}{section}{Литература к главе 3}
\markright{Литература к главе 3}

{\small
\paragraph{Использованная в исходном учебнике}
\begin{enumerate}[label=\arabic*.]
\item Жданов Г.\,С. Физика твёрдого тела.~— М.: МГУ, 1961.~— 501~с.
\item Федоров Г.\,Б., Яковлев Е.\,И. Металлофизика реакторных материалов.~—
М.: МИФИ, 1984.~— 76~с.
\item Федоров Г.\,Б., Яковлев Е.\,И. Атомное строение металлов и
сплавов.~— М.: МИФИ, 1985.~— 64~с.
\item Смирнов Е.\,А. Термодинамика фазовых превращений.~— М.: МИФИ,
1998.~— 83~с.
\item Нечаев В.\,В., Смирнов Е.\,А. Физическая химия сплавов.~— М.: МИФИ,
2006.~— 227~с.
\item Смирнов Е.\,А. Применение термодинамики необратимых явлений к анализу
диффузионных явлений в реакторных материалах.~— М.: МИФИ, 1984.~— 38~с.
\item Гуров К.\,П., Смирнов Е.\,А., Шабалин А.\,Н. Диффузия и кинетика
фазовых превращений в металлах и сплавах.~— М.: МИФИ, 1990.~— 80~с.
\item Даркен Л.\,С., Гурри Р.\,В. Физическая химия металлов.~— М.:
Металлургия, 1960.~— 582~с.
\item Лившиц Б.\,Г., Крапошин В.\,С., Линецкий Я.\,Л. Физические свойства
металлов и сплавов. 2-е изд.~— М.: Металлургия, 1980.~— 320~с.
\item Епифанов Г.\,И. Физика твёрдого тела. 2-е изд.~— М.: Высшая школа,
1977.~— 287~с.
\item Смирнов Е.\,А., Соколов Н.\,А., Яковлев Е.\,И. Физические свойства
металлов и сплавов.~— М.: МИФИ, 1992.~— 84~с.
\item Физические величины: справочник / под ред. И.\,С.~Григорьева,
Е.\,З.~Мейлихова.~— М.: Энергоатомиздат, 1991.~— 1232~с.
\end{enumerate}

\paragraph{Дополнительно}
\begin{enumerate}[label=\arabic*., resume]
\item Киттель Ч. Введение в физику твёрдого тела.~— М.: Наука, 1978.
\item Ашкрофт Н., Мермин Н. Физика твёрдого тела. В 2~т.~— М.: Мир, 1979.
\item Павлов П.\,В., Хохлов А.\,Ф. Физика твёрдого тела.~— М.: Высшая
школа, 2000.
\item Бокштейн Б.\,С. Диффузия в металлах.~— М.: Металлургия, 1978.
\item Новик А., Берри Б. Релаксационные явления в кристаллах.~— М.:
Атомиздат, 1975.
\item Вонсовский С.\,В. Магнетизм.~— М.: Наука, 1971.
\end{enumerate}
}
